\section{Ambiente de Desenvolvimento e Automação}

O fluxo de engenharia digital adotado combina ferramentas de síntese para lógica programável e utilitários abertos de simulação em ambiente Linux, estruturados em uma esteira de automação por linha de comando.

\subsection{Plataforma de Síntese e Dispositivo Alvo}
A síntese lógica e o mapeamento tecnológico foram configurados no software Intel Quartus Prime Standard/Lite Edition \cite{intel_quartus_handbook}. O dispositivo físico de destino é o FPGA Intel MAX 10 modelo \texttt{10M50DAF484C7G}, presente na placa de desenvolvimento Terasic DE10-Lite \cite{terasic_de10_lite_user_manual}. Este dispositivo possui 50.000 elementos lógicos, blocos de memória embarcada M9K de 432 Kibit e blocos multiplicadores dedicados de $18\times18$.

O processo de compilação física divide-se em três etapas acionadas via terminal:
\begin{enumerate}
    \item \texttt{quartus\_map}: análise sintática e síntese lógica em portas genéricas independentes de tecnologia;
    \item \texttt{quartus\_fit}: posicionamento e roteamento nos recursos de silício e canais de interconexão física do MAX 10;
    \item \texttt{quartus\_asm}: geração do arquivo binário de configuração final em formato \texttt{.sof}.
\end{enumerate}

A gravação física na placa é realizada via interface JTAG com o gravador USB-Blaster integrado, utilizando o utilitário \texttt{quartus\_pgm} após detecção com \texttt{jtagconfig}.

\subsection{Ferramentas de Simulação e Verificação RTL}
Para acelerar o ciclo de depuração temporal, empregou-se o compilador de simulação Icarus Verilog através do comando \texttt{iverilog} e seu motor de execução \texttt{vvp} \cite{icarus_verilog}. As formas de onda dos barramentos e sinais de controle foram exportadas para arquivos em formato VCD e inspecionadas graficamente no visualizador GTKWave.

\subsection{Automação via Makefile e Scripts de Regressão}
Todo o fluxo de compilação, simulação e síntese foi encapsulado em um \texttt{Makefile} centralizado e scripts Bash estruturados no diretório \texttt{scripts/}:
\begin{itemize}
    \item \texttt{make test}: executa o script \texttt{run\_all\_tests.sh}, disparando a regressão sobre as 7 bancadas com verificação de status;
    \item \texttt{make top}: simula o processador executando o código de máquina \texttt{sort.bin};
    \item \texttt{make quartus}: executa a análise estrutural rápida de síntese e elaboração;
    \item \texttt{make quartus-build}: conduz a compilação completa até a geração do bitstream \texttt{.sof}.
\end{itemize}
Essa automação garante reprodutibilidade estrita em ambientes de integração contínua, eliminando dependências manuais de telas gráficas.
